% ==============================================================================
% Abstract (structured per JQC requirements)
% ==============================================================================
\begin{abstract}
\noindent\textbf{Objectives.}
Prior spatial analyses of disappearances in Mexico treat them as a single phenomenon.
We present the first outcome-disaggregated spatial analysis, examining whether four outcome statuses---Not Located, Located Alive, Located Dead, and Total---produce distinct hotspot geographies.

\noindent\textbf{Methods.}
We analyze \nmunis{} municipalities over \nmonths{} months (January 2015--December 2025) using Local Indicators of Spatial Association (LISA) with Queen contiguity weights (\nLISAruns{} runs: 4 statuses $\times$ 3 sexes $\times$ \nmonths{} months).
Geographic concentration is measured via Gini coefficients benchmarked against the law of crime concentration.
Cross-status spatial overlap is quantified using Jaccard similarity on pairwise High-High cluster sets.
Trend slopes use Newey-West HAC standard errors.

\noindent\textbf{Results.}
Outcome statuses are spatially non-interchangeable: Jaccard similarity between Located Dead and Not Located hotspots is 0.143; between Located Dead and Located Alive, 0.151.
Concentration exceeds the Weisburd benchmark two- to threefold: 50\% of Not Located cases concentrate in 42 municipalities (1.7\%).
Centro is the fastest-growing cluster region (Not Located: $+4.96$ municipalities/year, $p < 0.001$), invisible to prior annual analyses.
The ``rising Baj\'{i}o'' hypothesis is rejected (Total: $-1.09$/year, $p = 0.024$).
Sex disaggregation reveals divergent male/female hotspot geographies for Not Located (Jaccard = 0.262) and a power failure for female Located Dead (zero HH municipalities).
Findings align with the UN Committee on Enforced Disappearances' 2026 characterization of spatially heterogeneous attacks.

\noindent\textbf{Conclusions.}
Forensic and prevention resources require different geographic targeting: Located Dead hotspots and Not Located hotspots are different places.
Composite disappearance maps conflating all statuses misallocate resources toward Located Alive--dominated geographies.
\end{abstract}
